\section{Conclusion and Discussion}
\label{sec:conclusion}

This paper presents Point Transformer V3, a stride towards overcoming the traditional trade-offs between accuracy and efficiency in point cloud processing. Guided by a novel interpretation of the \textbf{scaling principle} in backbone design, we propose that model performance is more profoundly influenced by scale than by complex design intricacies. By prioritizing efficiency over the accuracy of less impactful mechanisms, we harness the power of scale, leading to enhanced performance. Simply put, by making the model \textbf{simpler} and \textbf{faster}, we enable it to become \textbf{stronger}.

We discuss \textit{limitations and broader impacts} as follows:
\begin{itemize}[leftmargin=4mm, itemsep=0mm, topsep=0mm, partopsep=0mm]
    \item \textit{Attention mechanisum.} In prioritizing efficiency, PTv3 reverts to utilizing dot-product attention, which has been well-optimized through engineering efforts. However, we do note a reduction in convergence speed and a limitation in further scaling depth compared to vector attention. This issue also observed in recent advancements in transformer technology~\cite{xiao2023streamingllm}, is attributed to ``attention sinks'' stemming from the dot-product and softmax operations. Consequently, our findings reinforce the need for continued exploration of attention mechanisms.
    \item \textit{Scaling parameters.} PTv3 transcends the existing trade-offs between accuracy and efficiency, paving the way for investigating 3D transformers at larger parameter scales within given computational resources. While this exploration remains a topic for future work, current point cloud transformers already demonstrate an over-capacity for existing tasks. We advocate for a combined approach that scales up both the model parameters and the scope of data and tasks (e.g., learning from all available data, multi-task frameworks, and multi-modality tasks). Such an integrated strategy could fully unlock the potential of scaling in 3D representation learning.
    \item \textit{Multiple modalities.} Point cloud serialization provides a robust methodology for transforming n-dimensional data into a structured 1D format, effectively preserving spatial proximity. This technique can similarly be applied to image data, enabling its conversion into a language-style 1D structure that PTv3 can efficiently encode. This capability opens new avenues for the development of multimodal models that bridge 2D and 3D spaces, fostering opportunities for large-scale, synergistic pre-training that integrates both image and point cloud data.
\end{itemize}
